\documentclass[a4paper,10pt]{report}\input{../0.Préambule/Préambule_cours_prof}\begin{document}

\newpage
\section*{Applications de la Proportionnalité}
\addcontentsline{toc}{section}{Applications de la Proportionnalité}

\subsection*{Vitesse}
\addcontentsline{toc}{subsection}{Vitesse}

\begin{exer1}
\medskip

\begin{enumerate}[font=\color{exer1}]
\item  Le record du monde du $100$m est détenu par le jamaïcain Usain Bolt en $9,58$s. Quelle a été sa vitesse en m/s lors de sa course ?
\item Le record du monde du $\np{10 000}$m est détenu par l'ougandais Joshua Cheptegei en $26$min $11$s. Quelle a été sa vitesse en m/s puis en km/h lors de sa course ?
\end{enumerate}
\end{exer1}

\begin{exer1}
Cynthia est partie de chez elle à $8$h$30$ et est arrivée à son lieu de vacances à $16$h$50$ après avoir parcouru $625$km en voiture.

\medskip
\noindent Quelle a été la vitesse moyenne du trajet ?
\end{exer1}

\begin{exer1}
Une nuée ardente composée de gaz surchauffés, de cendre, de pierre ponce et de roche pulvérisée s'échappe latéralement à une vitesse initiale de $350$km/h et accélère rapidement pour atteindre les $1080$km/h.

\medskip
\noindent Quelle distance (en km) la nuée ardente a-t-elle parcourue en $30$s à sa vitesse maximale ?
\end{exer1}

\begin{exer2}
Le TGV « Nord » part de Lille à $10$h$20$ vers Paris à la vitesse de $227$km.h$^{-1}$ et le TGV « Sud » part de Paris à $10$h$30$ vers Lille à la vitesse de $239$ km.h$^{-1}$. La distance Lille-Paris est environ de $220$km par le train. 

\medskip
\noindent Ces deux trains vont-ils se croiser avant $10$h$53$ ?
\end{exer2}

\subsection*{Pourcentages}
\addcontentsline{toc}{subsection}{Pourcentages}

\begin{exer1}
Timéo veut emprunter $\np{3 000}$ euros. À quelle banque va-t-il s'adresser ?
\begin{center}
\renewcommand{\arraystretch}{2}
\begin{tabularx}{\linewidth}{|>{\columncolor{exer1!20}}p{3cm}|*{2}{>{\centering \arraybackslash}X|}}
\hline
Raison sociale & Banque du Nord & Banque du Sud \\\hline
Coût du crédit & $2,5\%$ du capital emprunté & $3,2\%$ du capital emprunté \\\hline
Assurance & $200$ euros & $155$ euros \\\hline
\end{tabularx}
\end{center}
\end{exer1}

\begin{exer1}
Le collège Victor Hugo accueille $800$ élèves dont $480$ garçons.
$30\%$ des garçons et $20\%$ des filles s'y rendent à pied.
\begin{enumerate}[font=\color{exer1}]
\item  Pour déterminer le pourcentage d'élèves qui vont à pied au collège, Chloé, Maréva et Teihotu discutent.

\medskip
Chloé dit : «  $20 +30 =50$ donc $50\%$ des élèves y vont à pied. »

\medskip
Maréva répond : « Tu as oublié de diviser par $2$. $25\%$ des élèves y vont à pied. »

\medskip
Teihotu rajoute : « On ne peut pas deviner le résultat car il n'y a pas le même nombre de garçons et de filles. »

\medskip
Commenter le raisonnement de chacun de ces élèves.

\item \begin{enumerate}[font=\color{exer1}]
\item  Calculer le nombre de garçons qui vont à pied à ce collège.
\item  Calculer le nombre de filles qui vont à pied à ce collège.
\item  En déduire le pourcentage d'élèves qui vont à pied au collège Victor Hugo.
\end{enumerate}
\end{enumerate}
\end{exer1}

\begin{exer1}
Un zoo abrite $\np{3 600}$ animaux dont $15\%$ représentent des espèces en voie de disparition. Les singes représentent $5\%$ des espèces en voie de disparition et $20\%$ des autres espèces.

\medskip
\noindent Calculer le pourcentage de singes dans ce zoo.
\end{exer1}

\begin{exer1}
Dans une classe, il y a $15$ élèves qui étudient l'allemand. Ces $15$ élèves représentent $60\%$ des élèves de la classe.

\medskip
\noindent Combien y-a-t-il d'élèves dans la classe ?
\end{exer1}

\begin{exer1}
Un congrès de scientifiques s'est divisé en deux commissions. Dans la première commission de$ 20$ personnes, il y a $15\%$ de femmes.
Dans la deuxième commission de $60$ personnes, il y a $25\%$ de femmes.

\medskip
\noindent Quel est le pourcentage de femmes dans ce congrès ?
\end{exer1}

\begin{exer1}
Dans le parking de la mairie, il y a $150$ voitures dont $80\%$ de marque française. Dans le parking du lac, il y a $250$ voitures dont $96\%$ de marque française.

\medskip
\noindent Quel est le pourcentage de voitures de marque française dans les deux parkings réunis ?
\end{exer1}

\begin{exer1}
Dans un troupeau de $120$ animaux, il y a $60\%$ de moutons. Dans un autre troupeau de $180$ animaux, il y a $40\%$ de moutons. Pour monter à l'alpage, les deux troupeaux sont rassemblés. 

\medskip
\noindent Quel est le pourcentage de moutons dans les deux groupes réunis ?
\end{exer1}

\begin{exer1}
Les Kangiens parlent soit seulement l'anglais, soit seulement le français, soit les deux langues. $85\%$ parlent anglais, $75\%$ parlent français.

\medskip 
\noindent Quel pourcentage de Kangiens est bilingue :
\begin{enumerate}[font=\color{exer1}]
\begin{tabularx}{\linewidth}{*{5}{>{\centering \arraybackslash}X}}
\item $50\%$ & \item $57\%$ & \item $25\%$ & \item $60\%$ & \item $40\%$ \\
\end{tabularx}
\end{enumerate}
\end{exer1}

\begin{exer2}
Un camion vide pèse $\np{2 000}$kg. Ce matin, la cargaison représentait $80\%$ du poids total du camion chargé. Au premier arrêt, on a déchargé le quart de la cargaison.
 
\medskip 
\noindent Quel pourcentage du poids total représente alors la cargaison :
\begin{enumerate}[font=\color{exer2}]
\begin{tabularx}{\linewidth}{*{5}{>{\centering \arraybackslash}X}}
\item $20\%$ & \item $25\%$ & \item $55\%$ & \item $60\%$ & \item $75\%$ \\
\end{tabularx}
\end{enumerate}
\end{exer2}

\begin{exer2}
Le cours de mathématiques a commencé à $10$h. Lucas regarde l'heure, il est $10$h$35$min. « Tiens, dit-il, il y a déjà $70\%$ du cours qui est passé ! ».

\medskip
\noindent Quelle heure est-il ?
\end{exer2}

\begin{exer3}
Lors de la crue de l'Ouvèze (affluent du Rhône) qui fit $42$ morts le $22$ septembre $1992$, on a estimé que le débit de cette rivière avait atteint un maximum de $\np{1 100}$m$^3$.s$^{-1}$ alors que le débit moyen est de $5,2$m$^3$.s$^{-1}$. Quel pourcentage d'augmentation cela représente-t-il ?
\end{exer3}

\end{document}